\chapter{Conclusion}

\section{Restatement of the Research Aim}

\begin{itemize}
    \item Restate the aim: to investigate whether compact recurrent neural networks can emulate the Boss DS-1 for real-time distortion plugin use.
    \item Restate the architectural comparison: GRU versus LSTM under matched data, training, and deployment conditions.
    \item Reconnect the work to the title of the dissertation and the practice-based research framing.
\end{itemize}

\section{Direct Answer to the Research Question}

\begin{itemize}
    \item Answer whether the models achieved an acceptable objective approximation of the DS-1 according to the final ESR and related metrics.
    \item Answer which recurrent architecture offered the better accuracy-to-efficiency trade-off.
    \item State clearly whether this conclusion is based on objective metrics, deployment observations, or both.
    \item Avoid overstating conclusions beyond the evidence collected.
\end{itemize}

\section{Summary of Contributions}

\begin{itemize}
    \item An original matched dry/wet DS-1 recording dataset at a documented fixed knob setting.
    \item A reproducible Python training workflow using compact recurrent neural models.
    \item A controlled comparison between LSTM and GRU architectures using the same device, data split, and evaluation framework.
    \item Exported model weights in a JSON format suitable for C\texttt{++} inference.
    \item A JUCE/RTNeural plugin artefact that demonstrates the deployment path from trained model to real-time audio software.
    \item Documentation of practical constraints, including real-time safety and the need for sample-rate conversion around the 44.1~kHz-trained model.
\end{itemize}

\section{Research Limitations}

\begin{itemize}
    \item One physical DS-1 unit was used, so the results do not represent every DS-1 unit or hardware revision.
    \item One knob setting was captured, so the model is not a parameterised emulation of Tone, Distortion, and Level.
    \item No formal listening study was conducted, so perceptual equivalence remains unproven.
    \item The plugin should be understood as a research artefact rather than a polished commercial release.
    \item If the sample-rate conversion path remains incomplete, non-44.1~kHz DAW operation must be treated as an unresolved deployment limitation.
\end{itemize}

\section{Future Work}

\begin{itemize}
    \item Train parameter-conditioned models that accept Tone, Distortion, and Level settings as inputs.
    \item Expand the dataset across multiple knob settings, playing styles, guitars, and gain levels.
    \item Validate across multiple physical DS-1 units to study unit-to-unit variation.
    \item Conduct a formal perceptual study using MUSHRA, ABX, or a comparable listening protocol with trained listeners.
    \item Implement and regression-test host sample-rate to 44.1~kHz to host sample-rate conversion for plugin inference.
    \item Add automated plugin tests and broader host validation across DAWs, plugin formats, buffer sizes, and sample rates.
    \item Profile runtime performance more systematically across model sizes and hardware platforms.
\end{itemize}

\section{Final Reflection}

\begin{itemize}
    \item Reflect on the value of the project as both a technical artefact and a research process.
    \item Discuss how the project demonstrates the practical connection between machine learning, digital signal processing, and music technology.
    \item Close by returning to the central motivation: preserving the expressive qualities of analog distortion in a modern real-time software workflow.
\end{itemize}
